% --- System overview: Uniform Herding configuration ---
% Place with:  % --- System overview: Uniform Herding configuration ---
% Place with:  % --- System overview: Uniform Herding configuration ---
% Place with:  % --- System overview: Uniform Herding configuration ---
% Place with:  \input{data/diagram/system_overview.tex}
% Preamble requires:
%   \usepackage{tikz}
%   \usetikzlibrary{arrows.meta, fit, backgrounds, positioning}

\definecolor{compfill}{RGB}{240,244,248}
\definecolor{compdraw}{RGB}{140,155,170}
\definecolor{ablfill}{RGB}{255,243,224}
\definecolor{abldraw}{RGB}{230,126,34}
\definecolor{arrowcol}{RGB}{90,90,90}
\definecolor{notecol}{RGB}{120,120,120}
\definecolor{regionfill}{RGB}{250,250,250}
\definecolor{regiondraw}{RGB}{210,210,210}
\definecolor{labelcol}{RGB}{150,150,150}
\definecolor{mergecol}{RGB}{160,170,180}
\begin{figure}[t]
\centering
\begin{tikzpicture}[
  scale=0.9,
  >=Stealth,
  node distance=4pt,
  base/.style={
    draw, rounded corners=2pt, minimum height=8.5mm,
    minimum width=24mm, align=center, inner sep=3pt,
    font=\sffamily\small, line width=0.5pt
  },
  comp/.style={base, fill=compfill, draw=compdraw},
  abl/.style={base, fill=ablfill, draw=abldraw, line width=0.7pt},
  flow/.style={->, line width=0.55pt, draw=arrowcol},
  back/.style={flow, dashed},
  note/.style={font=\sffamily\tiny, text=notecol},
  region/.style={
    draw=regiondraw, fill=regionfill, rounded corners=4pt,
    inner sep=9pt, line width=0.3pt
  },
  rlabel/.style={font=\sffamily\scriptsize\bfseries, text=labelcol},
]

% ==================== TRAINING ====================
\node[comp] (data) at (0, 0)       {Task $t$ data};
\node[comp] (rpl)  at (0, -1.3)    {Replay};
\node[circle, fill=mergecol, inner sep=1.5pt]
             (mrg) at (3, -0.65)   {};
\node[comp] (bb)   at (6, -0.65)   {Backbone $f_\theta$};
\node[abl]  (head) at (9, -0.65)   {Head};
\node[abl]  (loss) at (12, -0.65)  {Loss};
\node[comp] (tch)  at (12, -2.1)   {Teacher $f_{\theta'}$};

\draw[flow] (data.east) -- (mrg);
\draw[flow] (rpl.east)  -- (mrg);
\draw[flow] (mrg)  -- (bb);
\draw[flow] (bb)   -- (head);
\draw[flow] (head) -- (loss);
\draw[flow] (tch)  -- (loss);

\node[note, below=2pt of head] {cosine-margin\,/\,linear};
\node[note, above=2pt of loss] {CE + KD\,/\,CE only};
\node[note, below=2pt of rpl]  {$b$ samples};

% headroom reserved above the training row, for the panel title
\coordinate (R1pad) at ($(data.north)+(0,0.5)$);

% ==================== MEMORY REBUILD ====================
\node[comp] (pool) at (1, -4.7)    {Class pools};
\node[abl]  (sel)  at (4.8, -4.7)  {Selection};
\node[comp] (mem)  at (8.6, -4.7)  {Memory $\mathcal{M}$};
\node[comp] (prt)  at (12, -4.7)   {Prototypes $\hat{\mu}_c$};

\draw[flow] (pool) -- (sel);
\draw[flow] (sel)  -- (mem);
\draw[flow] (mem)  -- (prt);

\node[note, below=2pt of sel] {herding\,/\,random};
\node[note, above=2pt of sel] {\itshape rebuild in current $f_\theta$};
\node[note, below=2pt of mem] {budget $M$};

% headroom above/below the memory-rebuild row
\coordinate (R2padT) at ($(pool.north)+(0,0.5)$);
\coordinate (R2padB) at ($(mem.south)+(0,-0.4)$);

% ==================== EVALUATION ====================
\node[comp] (qry)  at (1, -7.74)   {Query $x$};
\node[comp] (bb2)  at (4.8, -7.74) {Backbone $f_\theta$};
\node[abl]  (rdo)  at (8.6, -7.74) {Readout};
\node[comp] (pred) at (12, -7.74)  {$\hat{y}$};

\draw[flow] (qry) -- (bb2);
\draw[flow] (bb2) -- (rdo);
\draw[flow] (rdo) -- (pred);

\node[note, below=2pt of rdo] {NME\,/\,head-logit};

% headroom above/below the evaluation row
\coordinate (R3padT) at ($(qry.north)+(0,0.5)$);
\coordinate (R3padB) at ($(rdo.south)+(0,-0.4)$);

% ==================== INTER-REGION CONNECTIONS ====================
% Memory feeds replay samples -- jog routed through the neutral band
% between the two panel borders, offset entry point clears "b samples"
\draw[back] (mem.north) -- ++(0, 1.13) -| ($(rpl.south)+(-1.0,0)$);
% Prototypes supply the readout -- jog routed through the neutral band
% between the two panel borders
\draw[back] (prt.south) -- ++(0, -1.05) -| (rdo.north);

% ==================== REGION BACKGROUNDS ====================
\begin{scope}[on background layer]
  \node[region, fit=(data)(rpl)(loss)(tch)(R1pad)]            (R1) {};
  \node[region, fit=(pool)(prt)(R2padT)(R2padB)]               (R2) {};
  \node[region, fit=(qry)(pred)(R3padT)(R3padB)]                (R3) {};
\end{scope}

\node[rlabel, anchor=north west]
  at ([shift={(5pt,-5pt)}]R1.north west)
  {Training (task $t \geq 1$)};
\node[rlabel, anchor=north west]
  at ([shift={(5pt,-5pt)}]R2.north west)
  {Memory rebuild (after each task)};
\node[rlabel, anchor=north west]
  at ([shift={(5pt,-5pt)}]R3.north west)
  {Evaluation};

\end{tikzpicture}

\caption{Overview of the Uniform Herding replay configuration.  Orange boxes mark
the four components varied in the ablation study: classifier head, training
loss, exemplar selection, and evaluation readout; the Uniform Herding choice
is listed first in each annotation.  Dashed arrows show the memory feedback loop
and prototype supply.  The resource sweeps additionally vary the memory
budget~$M$ and the per-minibatch retrieval count~$b$.}
\label{fig:system-overview}
\end{figure}

% Preamble requires:
%   \usepackage{tikz}
%   \usetikzlibrary{arrows.meta, fit, backgrounds, positioning}

\definecolor{compfill}{RGB}{240,244,248}
\definecolor{compdraw}{RGB}{140,155,170}
\definecolor{ablfill}{RGB}{255,243,224}
\definecolor{abldraw}{RGB}{230,126,34}
\definecolor{arrowcol}{RGB}{90,90,90}
\definecolor{notecol}{RGB}{120,120,120}
\definecolor{regionfill}{RGB}{250,250,250}
\definecolor{regiondraw}{RGB}{210,210,210}
\definecolor{labelcol}{RGB}{150,150,150}
\definecolor{mergecol}{RGB}{160,170,180}
\begin{figure}[t]
\centering
\begin{tikzpicture}[
  scale=0.9,
  >=Stealth,
  node distance=4pt,
  base/.style={
    draw, rounded corners=2pt, minimum height=8.5mm,
    minimum width=24mm, align=center, inner sep=3pt,
    font=\sffamily\small, line width=0.5pt
  },
  comp/.style={base, fill=compfill, draw=compdraw},
  abl/.style={base, fill=ablfill, draw=abldraw, line width=0.7pt},
  flow/.style={->, line width=0.55pt, draw=arrowcol},
  back/.style={flow, dashed},
  note/.style={font=\sffamily\tiny, text=notecol},
  region/.style={
    draw=regiondraw, fill=regionfill, rounded corners=4pt,
    inner sep=9pt, line width=0.3pt
  },
  rlabel/.style={font=\sffamily\scriptsize\bfseries, text=labelcol},
]

% ==================== TRAINING ====================
\node[comp] (data) at (0, 0)       {Task $t$ data};
\node[comp] (rpl)  at (0, -1.3)    {Replay};
\node[circle, fill=mergecol, inner sep=1.5pt]
             (mrg) at (3, -0.65)   {};
\node[comp] (bb)   at (6, -0.65)   {Backbone $f_\theta$};
\node[abl]  (head) at (9, -0.65)   {Head};
\node[abl]  (loss) at (12, -0.65)  {Loss};
\node[comp] (tch)  at (12, -2.1)   {Teacher $f_{\theta'}$};

\draw[flow] (data.east) -- (mrg);
\draw[flow] (rpl.east)  -- (mrg);
\draw[flow] (mrg)  -- (bb);
\draw[flow] (bb)   -- (head);
\draw[flow] (head) -- (loss);
\draw[flow] (tch)  -- (loss);

\node[note, below=2pt of head] {cosine-margin\,/\,linear};
\node[note, above=2pt of loss] {CE + KD\,/\,CE only};
\node[note, below=2pt of rpl]  {$b$ samples};

% headroom reserved above the training row, for the panel title
\coordinate (R1pad) at ($(data.north)+(0,0.5)$);

% ==================== MEMORY REBUILD ====================
\node[comp] (pool) at (1, -4.7)    {Class pools};
\node[abl]  (sel)  at (4.8, -4.7)  {Selection};
\node[comp] (mem)  at (8.6, -4.7)  {Memory $\mathcal{M}$};
\node[comp] (prt)  at (12, -4.7)   {Prototypes $\hat{\mu}_c$};

\draw[flow] (pool) -- (sel);
\draw[flow] (sel)  -- (mem);
\draw[flow] (mem)  -- (prt);

\node[note, below=2pt of sel] {herding\,/\,random};
\node[note, above=2pt of sel] {\itshape rebuild in current $f_\theta$};
\node[note, below=2pt of mem] {budget $M$};

% headroom above/below the memory-rebuild row
\coordinate (R2padT) at ($(pool.north)+(0,0.5)$);
\coordinate (R2padB) at ($(mem.south)+(0,-0.4)$);

% ==================== EVALUATION ====================
\node[comp] (qry)  at (1, -7.74)   {Query $x$};
\node[comp] (bb2)  at (4.8, -7.74) {Backbone $f_\theta$};
\node[abl]  (rdo)  at (8.6, -7.74) {Readout};
\node[comp] (pred) at (12, -7.74)  {$\hat{y}$};

\draw[flow] (qry) -- (bb2);
\draw[flow] (bb2) -- (rdo);
\draw[flow] (rdo) -- (pred);

\node[note, below=2pt of rdo] {NME\,/\,head-logit};

% headroom above/below the evaluation row
\coordinate (R3padT) at ($(qry.north)+(0,0.5)$);
\coordinate (R3padB) at ($(rdo.south)+(0,-0.4)$);

% ==================== INTER-REGION CONNECTIONS ====================
% Memory feeds replay samples -- jog routed through the neutral band
% between the two panel borders, offset entry point clears "b samples"
\draw[back] (mem.north) -- ++(0, 1.13) -| ($(rpl.south)+(-1.0,0)$);
% Prototypes supply the readout -- jog routed through the neutral band
% between the two panel borders
\draw[back] (prt.south) -- ++(0, -1.05) -| (rdo.north);

% ==================== REGION BACKGROUNDS ====================
\begin{scope}[on background layer]
  \node[region, fit=(data)(rpl)(loss)(tch)(R1pad)]            (R1) {};
  \node[region, fit=(pool)(prt)(R2padT)(R2padB)]               (R2) {};
  \node[region, fit=(qry)(pred)(R3padT)(R3padB)]                (R3) {};
\end{scope}

\node[rlabel, anchor=north west]
  at ([shift={(5pt,-5pt)}]R1.north west)
  {Training (task $t \geq 1$)};
\node[rlabel, anchor=north west]
  at ([shift={(5pt,-5pt)}]R2.north west)
  {Memory rebuild (after each task)};
\node[rlabel, anchor=north west]
  at ([shift={(5pt,-5pt)}]R3.north west)
  {Evaluation};

\end{tikzpicture}

\caption{Overview of the Uniform Herding replay configuration.  Orange boxes mark
the four components varied in the ablation study: classifier head, training
loss, exemplar selection, and evaluation readout; the Uniform Herding choice
is listed first in each annotation.  Dashed arrows show the memory feedback loop
and prototype supply.  The resource sweeps additionally vary the memory
budget~$M$ and the per-minibatch retrieval count~$b$.}
\label{fig:system-overview}
\end{figure}

% Preamble requires:
%   \usepackage{tikz}
%   \usetikzlibrary{arrows.meta, fit, backgrounds, positioning}

\definecolor{compfill}{RGB}{240,244,248}
\definecolor{compdraw}{RGB}{140,155,170}
\definecolor{ablfill}{RGB}{255,243,224}
\definecolor{abldraw}{RGB}{230,126,34}
\definecolor{arrowcol}{RGB}{90,90,90}
\definecolor{notecol}{RGB}{120,120,120}
\definecolor{regionfill}{RGB}{250,250,250}
\definecolor{regiondraw}{RGB}{210,210,210}
\definecolor{labelcol}{RGB}{150,150,150}
\definecolor{mergecol}{RGB}{160,170,180}
\begin{figure}[t]
\centering
\begin{tikzpicture}[
  scale=0.9,
  >=Stealth,
  node distance=4pt,
  base/.style={
    draw, rounded corners=2pt, minimum height=8.5mm,
    minimum width=24mm, align=center, inner sep=3pt,
    font=\sffamily\small, line width=0.5pt
  },
  comp/.style={base, fill=compfill, draw=compdraw},
  abl/.style={base, fill=ablfill, draw=abldraw, line width=0.7pt},
  flow/.style={->, line width=0.55pt, draw=arrowcol},
  back/.style={flow, dashed},
  note/.style={font=\sffamily\tiny, text=notecol},
  region/.style={
    draw=regiondraw, fill=regionfill, rounded corners=4pt,
    inner sep=9pt, line width=0.3pt
  },
  rlabel/.style={font=\sffamily\scriptsize\bfseries, text=labelcol},
]

% ==================== TRAINING ====================
\node[comp] (data) at (0, 0)       {Task $t$ data};
\node[comp] (rpl)  at (0, -1.3)    {Replay};
\node[circle, fill=mergecol, inner sep=1.5pt]
             (mrg) at (3, -0.65)   {};
\node[comp] (bb)   at (6, -0.65)   {Backbone $f_\theta$};
\node[abl]  (head) at (9, -0.65)   {Head};
\node[abl]  (loss) at (12, -0.65)  {Loss};
\node[comp] (tch)  at (12, -2.1)   {Teacher $f_{\theta'}$};

\draw[flow] (data.east) -- (mrg);
\draw[flow] (rpl.east)  -- (mrg);
\draw[flow] (mrg)  -- (bb);
\draw[flow] (bb)   -- (head);
\draw[flow] (head) -- (loss);
\draw[flow] (tch)  -- (loss);

\node[note, below=2pt of head] {cosine-margin\,/\,linear};
\node[note, above=2pt of loss] {CE + KD\,/\,CE only};
\node[note, below=2pt of rpl]  {$b$ samples};

% headroom reserved above the training row, for the panel title
\coordinate (R1pad) at ($(data.north)+(0,0.5)$);

% ==================== MEMORY REBUILD ====================
\node[comp] (pool) at (1, -4.7)    {Class pools};
\node[abl]  (sel)  at (4.8, -4.7)  {Selection};
\node[comp] (mem)  at (8.6, -4.7)  {Memory $\mathcal{M}$};
\node[comp] (prt)  at (12, -4.7)   {Prototypes $\hat{\mu}_c$};

\draw[flow] (pool) -- (sel);
\draw[flow] (sel)  -- (mem);
\draw[flow] (mem)  -- (prt);

\node[note, below=2pt of sel] {herding\,/\,random};
\node[note, above=2pt of sel] {\itshape rebuild in current $f_\theta$};
\node[note, below=2pt of mem] {budget $M$};

% headroom above/below the memory-rebuild row
\coordinate (R2padT) at ($(pool.north)+(0,0.5)$);
\coordinate (R2padB) at ($(mem.south)+(0,-0.4)$);

% ==================== EVALUATION ====================
\node[comp] (qry)  at (1, -7.74)   {Query $x$};
\node[comp] (bb2)  at (4.8, -7.74) {Backbone $f_\theta$};
\node[abl]  (rdo)  at (8.6, -7.74) {Readout};
\node[comp] (pred) at (12, -7.74)  {$\hat{y}$};

\draw[flow] (qry) -- (bb2);
\draw[flow] (bb2) -- (rdo);
\draw[flow] (rdo) -- (pred);

\node[note, below=2pt of rdo] {NME\,/\,head-logit};

% headroom above/below the evaluation row
\coordinate (R3padT) at ($(qry.north)+(0,0.5)$);
\coordinate (R3padB) at ($(rdo.south)+(0,-0.4)$);

% ==================== INTER-REGION CONNECTIONS ====================
% Memory feeds replay samples -- jog routed through the neutral band
% between the two panel borders, offset entry point clears "b samples"
\draw[back] (mem.north) -- ++(0, 1.13) -| ($(rpl.south)+(-1.0,0)$);
% Prototypes supply the readout -- jog routed through the neutral band
% between the two panel borders
\draw[back] (prt.south) -- ++(0, -1.05) -| (rdo.north);

% ==================== REGION BACKGROUNDS ====================
\begin{scope}[on background layer]
  \node[region, fit=(data)(rpl)(loss)(tch)(R1pad)]            (R1) {};
  \node[region, fit=(pool)(prt)(R2padT)(R2padB)]               (R2) {};
  \node[region, fit=(qry)(pred)(R3padT)(R3padB)]                (R3) {};
\end{scope}

\node[rlabel, anchor=north west]
  at ([shift={(5pt,-5pt)}]R1.north west)
  {Training (task $t \geq 1$)};
\node[rlabel, anchor=north west]
  at ([shift={(5pt,-5pt)}]R2.north west)
  {Memory rebuild (after each task)};
\node[rlabel, anchor=north west]
  at ([shift={(5pt,-5pt)}]R3.north west)
  {Evaluation};

\end{tikzpicture}

\caption{Overview of the Uniform Herding replay configuration.  Orange boxes mark
the four components varied in the ablation study: classifier head, training
loss, exemplar selection, and evaluation readout; the Uniform Herding choice
is listed first in each annotation.  Dashed arrows show the memory feedback loop
and prototype supply.  The resource sweeps additionally vary the memory
budget~$M$ and the per-minibatch retrieval count~$b$.}
\label{fig:system-overview}
\end{figure}

% Preamble requires:
%   \usepackage{tikz}
%   \usetikzlibrary{arrows.meta, fit, backgrounds, positioning}

\definecolor{compfill}{RGB}{240,244,248}
\definecolor{compdraw}{RGB}{140,155,170}
\definecolor{ablfill}{RGB}{255,243,224}
\definecolor{abldraw}{RGB}{230,126,34}
\definecolor{arrowcol}{RGB}{90,90,90}
\definecolor{notecol}{RGB}{120,120,120}
\definecolor{regionfill}{RGB}{250,250,250}
\definecolor{regiondraw}{RGB}{210,210,210}
\definecolor{labelcol}{RGB}{150,150,150}
\definecolor{mergecol}{RGB}{160,170,180}
\begin{figure}[t]
\centering
\begin{tikzpicture}[
  scale=0.9,
  >=Stealth,
  node distance=4pt,
  base/.style={
    draw, rounded corners=2pt, minimum height=8.5mm,
    minimum width=24mm, align=center, inner sep=3pt,
    font=\sffamily\small, line width=0.5pt
  },
  comp/.style={base, fill=compfill, draw=compdraw},
  abl/.style={base, fill=ablfill, draw=abldraw, line width=0.7pt},
  flow/.style={->, line width=0.55pt, draw=arrowcol},
  back/.style={flow, dashed},
  note/.style={font=\sffamily\tiny, text=notecol},
  region/.style={
    draw=regiondraw, fill=regionfill, rounded corners=4pt,
    inner sep=9pt, line width=0.3pt
  },
  rlabel/.style={font=\sffamily\scriptsize\bfseries, text=labelcol},
]

% ==================== TRAINING ====================
\node[comp] (data) at (0, 0)       {Task $t$ data};
\node[comp] (rpl)  at (0, -1.3)    {Replay};
\node[circle, fill=mergecol, inner sep=1.5pt]
             (mrg) at (3, -0.65)   {};
\node[comp] (bb)   at (6, -0.65)   {Backbone $f_\theta$};
\node[abl]  (head) at (9, -0.65)   {Head};
\node[abl]  (loss) at (12, -0.65)  {Loss};
\node[comp] (tch)  at (12, -2.1)   {Teacher $f_{\theta'}$};

\draw[flow] (data.east) -- (mrg);
\draw[flow] (rpl.east)  -- (mrg);
\draw[flow] (mrg)  -- (bb);
\draw[flow] (bb)   -- (head);
\draw[flow] (head) -- (loss);
\draw[flow] (tch)  -- (loss);

\node[note, below=2pt of head] {cosine-margin\,/\,linear};
\node[note, above=2pt of loss] {CE + KD\,/\,CE only};
\node[note, below=2pt of rpl]  {$b$ samples};

% headroom reserved above the training row, for the panel title
\coordinate (R1pad) at ($(data.north)+(0,0.5)$);

% ==================== MEMORY REBUILD ====================
\node[comp] (pool) at (1, -4.7)    {Class pools};
\node[abl]  (sel)  at (4.8, -4.7)  {Selection};
\node[comp] (mem)  at (8.6, -4.7)  {Memory $\mathcal{M}$};
\node[comp] (prt)  at (12, -4.7)   {Prototypes $\hat{\mu}_c$};

\draw[flow] (pool) -- (sel);
\draw[flow] (sel)  -- (mem);
\draw[flow] (mem)  -- (prt);

\node[note, below=2pt of sel] {herding\,/\,random};
\node[note, above=2pt of sel] {\itshape rebuild in current $f_\theta$};
\node[note, below=2pt of mem] {budget $M$};

% headroom above/below the memory-rebuild row
\coordinate (R2padT) at ($(pool.north)+(0,0.5)$);
\coordinate (R2padB) at ($(mem.south)+(0,-0.4)$);

% ==================== EVALUATION ====================
\node[comp] (qry)  at (1, -7.74)   {Query $x$};
\node[comp] (bb2)  at (4.8, -7.74) {Backbone $f_\theta$};
\node[abl]  (rdo)  at (8.6, -7.74) {Readout};
\node[comp] (pred) at (12, -7.74)  {$\hat{y}$};

\draw[flow] (qry) -- (bb2);
\draw[flow] (bb2) -- (rdo);
\draw[flow] (rdo) -- (pred);

\node[note, below=2pt of rdo] {NME\,/\,head-logit};

% headroom above/below the evaluation row
\coordinate (R3padT) at ($(qry.north)+(0,0.5)$);
\coordinate (R3padB) at ($(rdo.south)+(0,-0.4)$);

% ==================== INTER-REGION CONNECTIONS ====================
% Memory feeds replay samples -- jog routed through the neutral band
% between the two panel borders, offset entry point clears "b samples"
\draw[back] (mem.north) -- ++(0, 1.13) -| ($(rpl.south)+(-1.0,0)$);
% Prototypes supply the readout -- jog routed through the neutral band
% between the two panel borders
\draw[back] (prt.south) -- ++(0, -1.05) -| (rdo.north);

% ==================== REGION BACKGROUNDS ====================
\begin{scope}[on background layer]
  \node[region, fit=(data)(rpl)(loss)(tch)(R1pad)]            (R1) {};
  \node[region, fit=(pool)(prt)(R2padT)(R2padB)]               (R2) {};
  \node[region, fit=(qry)(pred)(R3padT)(R3padB)]                (R3) {};
\end{scope}

\node[rlabel, anchor=north west]
  at ([shift={(5pt,-5pt)}]R1.north west)
  {Training (task $t \geq 1$)};
\node[rlabel, anchor=north west]
  at ([shift={(5pt,-5pt)}]R2.north west)
  {Memory rebuild (after each task)};
\node[rlabel, anchor=north west]
  at ([shift={(5pt,-5pt)}]R3.north west)
  {Evaluation};

\end{tikzpicture}

\caption{Overview of the Uniform Herding replay configuration.  Orange boxes mark
the four components varied in the ablation study: classifier head, training
loss, exemplar selection, and evaluation readout; the Uniform Herding choice
is listed first in each annotation.  Dashed arrows show the memory feedback loop
and prototype supply.  The resource sweeps additionally vary the memory
budget~$M$ and the per-minibatch retrieval count~$b$.}
\label{fig:system-overview}
\end{figure}
